\documentclass[10pt]{article}
\usepackage[T1]{fontenc}
\usepackage[utf8]{inputenc}
\usepackage[french]{babel}
\usepackage{amsmath,amssymb,mathrsfs,mathpazo}
\usepackage{geometry}
\geometry{a4paper,margin=2cm}
\usepackage{graphicx,array,tabularx,tabularray,colortbl,multirow,multicol,booktabs,enumitem,pifont,ulem,fancybox,diagbox,cancel}
\usepackage{pstricks,pstricks-add,pst-plot,pst-tree,pst-node,pst-text,pst-func,pst-math,pst-bspline,pst-3dplot,pst-eucl,pst-solides3d}
\let\pscalculate\undefined
\usepackage{pst-bezier}
\usepackage{tikz,pgfplots,tkz-tab}
\usetikzlibrary{arrows,arrows.meta,patterns,calc,positioning,trees,decorations.markings,shapes.geometric,shapes.misc,shapes.arrows,intersections,decorations.pathmorphing,decorations.pathreplacing,backgrounds,shadows,babel}
\pgfplotsset{compat=1.18}
\usepgfplotslibrary{fillbetween}
\usepackage[np]{numprint}
\usepackage{listings,makecell,pas-tableur,twoopt}
\let\barre\undefined
\usepackage[pstricks]{tablvar}
\definecolor{darkgreen}{RGB}{0,100,0}
\providecommand{\red}{\color{red}}
\providecommand{\blue}{\color{blue}}
\providecommand{\green}{\color{green}}
\providecommand{\black}{\color{black}}
\providecommand{\white}{\color{white}}
\providecommand{\gray}{\color{gray}}
\providecommand{\cyan}{\color{cyan}}
\providecommand{\magenta}{\color{magenta}}
\providecommand{\yellow}{\color{yellow}}
\providecommand{\violet}{\color{violet}}
\providecommand{\orange}{\color{orange}}
\providecommand{\euro}{\texteuro}
\providecommand{\EUR}{\texteuro}
\providecommand{\ieme}{\textsuperscript{e}}
\providecommand{\degres}{\ensuremath{^\circ}}
\providecommand{\degree}{\ensuremath{^\circ}}
\providecommand{\vv}[1]{\overrightarrow{#1}}
\providecommand{\Coord}[1]{#1}
\providecommand{\decosix}{}
\providecommand{\starredbullet}{\ensuremath{\bullet}}
\pagestyle{empty}
\setlength{\parindent}{0pt}
\usepackage{adjustbox}\DeclareUnicodeCharacter{2236}{:}\DeclareUnicodeCharacter{2190}{\ensuremath{\leftarrow}}
\providecommand{\R}{\mathbb{R}}
\providecommand{\N}{\mathbb{N}}
\providecommand{\Z}{\mathbb{Z}}
\providecommand{\Q}{\mathbb{Q}}
\providecommand{\D}{\mathbb{D}}
\providecommand{\C}{\mathbb{C}}
\providecommand{\e}{{\mathrm{e}}}
\providecommand{\vect}[1]{\overrightarrow{#1}}
\providecommand{\vectt}[1]{\overrightarrow{\mathrm{#1}}}
\providecommand{\vvt}[1]{\overrightarrow{\mathrm{#1}}}
\providecommand{\np}[2][]{{#2~#1}}
\providecommand{\vv}[1]{{\overrightarrow{#1}}}
\renewcommand{\d}{{\mathrm{d}}}
\providecommand{\degres}{{^\circ}}
\providecommand{\degre}{{^\circ}}
\providecommand{\blue}{\color{blue}}
\providecommand{\red}{\color{red}}
\providecommand{\enskip}{\quad}
\providecommand{\llbracket}{\lbrack\!\lbrack}
\providecommand{\rrbracket}{\rbrack\!\rbrack}
\providecommand{\euro}{{\text{€}}}
\providecommand{\no}{{\text{n°}}}
\providecommand{\interval}[2]{{[#1;#2]}}
\providecommand{\intervFO}[2]{{[#1;#2[}}
\providecommand{\dotfill}{\ldots}

\providecommand{\ds}{\displaystyle}
\providecommand{\textding}[1]{#1}
\providecommand{\Oij}{\(\mathrm{O};\vec{\imath},\vec{\jmath}\)}
\providecommand{\Oijk}{\(\mathrm{O};\vec{\imath},\vec{\jmath},\vec{k}\)}
\providecommand{\Ouv}{\(\mathrm{O};\vec{u},\vec{v}\)}
\providecommand{\intervFF}[2]{\ensuremath{[#1;#2]}}
\renewcommand{\euro}{\ifmmode\text{\texteuro}\else\texteuro\fi}
\DeclareUnicodeCharacter{00B0}{\ifmmode^{\circ}\else\textdegree\fi}
\DeclareUnicodeCharacter{20AC}{\ifmmode\text{\texteuro}\else\texteuro\fi}
% Énoncés : Éducation nationale. Transcription et corrigés : APMEP, Denis Vergès et contributeurs. Découpage et classement : CoopMaths, Sébastien Lozano, Éric Elter et contributeurs. Adaptation de la présentation : Maths Explorer.
% CC BY-SA : https://coopmaths.fr/www/about/
% https://coopmaths.fr/alea/static/e3c/2020/tex/e3c_2020_00_sujet4_4_cor.tex
\begin{document}
\noindent\textit{APMEP / CoopMaths. CC BY-SA. Présentation adaptée par Maths Explorer.}\par\medskip
\section*{Exercice 4 (5 points)}
	

	
	\subsection*{1. Calculer $y$ lorsque $x = 20$ cm}
	
	Le volume de la boîte est donc :
	\[
	V = 16xy
	\]
	Avec $V = 10000$ et $x = 20$, on a donc :
	\[
	10000 = 16 \times 20 \times y \iff 10000 = 320y \iff y = 31,25 \text{ cm}
	\]
	
	\subsection*{2. Pour toute valeur de $x > 0$, on note $f(x)$ l’aire du parallélépipède rectangle. Démontrer que : pour tout $x > 0$,}
	\[
	f(x) = \dfrac{20000}{x} + 32x + 625
	\]
	
	On a une base d’aire $xy$, deux côtés d’aire $16x$ et deux côtés d’aire $16y$.
	
	L’aire du parallélépipède rectangle est donc égale à :
	\[
	f(x) = xy + 32x + 32y \text{ et on sait que } 10000 = 16xy \iff y = \dfrac{625}{x}
	\]
	Donc
	\[
	f(x) = x \times \dfrac{625}{x} + 32x + 32 \times \dfrac{625}{x} = 625 + 32x + \dfrac{20000}{x} = \dfrac{20000}{x} + 32x + 625
	\]
	
	\subsection*{3. Quelles dimensions doit-on donner à ces boîtes pour que leur surface ait une aire minimale ?}
	
	On a $f'(x) = -\dfrac{20000}{x^2} + 32$.

	\begin{itemize}
		\item $\dfrac{20000}{x^2} + 32 > 0  \iff 32 > \dfrac{20000}{x^2} \iff x^2 > \dfrac{20000}{32} \iff x^2 > 625 \iff x > 25$
		\item $\dfrac{20000}{x^2} + 32 < 0  \iff 32 < \dfrac{20000}{x^2} \iff x^2 < \dfrac{20000}{32} \iff x^2 < 625 \iff x < 25$
		\item $\dfrac{20000}{x^2} + 32 = 0  \iff 32 = \dfrac{20000}{x^2} \iff x^2 = \dfrac{20000}{32} \iff x^2 = 625 \iff x = 25$
	\end{itemize}

	Donc l’aire est décroissante sur $[0, 25]$, puis croissante pour $x > 25$ : le minimum de l’aire est donc obtenu pour $x = 25$. On aura $y = \dfrac{625}{25} = 25$.
	
	Les boîtes feront donc : $25 \times 25 \times 16$ cm.

\end{document}
